% Pursuing Stacks (cotes 134-2 to 134-8): archive cote 134-4, site page 94 (PDF page 95), typescript page 181.
% typescript. Draft compared: ps3.tex:4819-4877.
% AI-transcribed by Claude Opus 5.5 (claude-opus-5-5, adaptive thinking, effort high), Message Batches API, batch ps2 (msgbatch_01NVwr1Yuo8qJf5rykeGBXPK), 2026-11-12; not yet checked by eye.
\scanpage{134-4/94}
\mspage{181}

(pages 159-163),  and notably the use of the word ``aspheric'' in the gene-%
ralization ``W-aspheric'' map (in a category M endowed with a saturated
set of arrows W)  -  which then practically obliged me,\xout{\illegible{1}} when
working in (Cat),  to call ``weakly aspheric'' a functor $C'\hand{\to}C$  which \xout{\illegible{1}}
\xout{\illegible{1}} spontaneously I surely would like to call simply ``aspheric''  -
and as a matter of fact,  it turned out I could'nt force myself to add a
``weakly'' before as I decreed I should  -  or if I did,  it was against a
very strong feeling of inapropriateness.  That decree \inserted{precisely} is an excellent il-%
lustration of loosing view of where the main emphasis belongs,  which I
would like now to make very clear to myself.

In all this work the underlying motivation or inspiration is geometri%
% NOTE: the right edge of the scan cuts the end of this line and of some following lines
co-topological,  and expressed technically quite accurately by the notion
of a topos and of maps (or ``morphisms'')of topoi,  and the wealth of geo-%
metric and algebraic intuitions \xout{\illegible{1}} which have developped around these.
One main point here is that topoi may be viewed as \uline{the} natural common
generalization of both topological spaces (the conventional support for
so-called topological intuition),  and of \xout{\illegible{1}} (small) categories,
\deleted{which} \inserted{where the latter} may be viewed as the ideal purely algebraic objects carrying topolo-%
gical information,  including all the conventional homology and homotopy
invariants.  This being so, the main emphasis in choice of terminology
in a context where working with small categories as ``spaces'', should
surely be in stressing throughout\hand{,}through the very wording\hand{,}the essential
identity between situations involving categories,  and corresponding situ%
ations involving topological spaces or topoi.  Thus It has been about
twenty years now that the \xout{\unsure{technical}} needs for developing etale cohomology
have told us a handfull of basic asphericity and acyclicity properties
for a map of topoi (which apparently have not yet been assimilated by
topologists\hand{,}in the context of \inserted[typed]{maps of} topological spaces ...),  including the
condition for such a map to be aspheric. This was recalled earlier
(page 37),  and the corresponding notion for a functor $f\colon C'\hand{\to}C$  was
introduced.  The name \uline{``aspheric map''} of topoi,\inserted{or} of topological spaces,
or of categories,  is \inserted{here} a perfectly suggestive one.  \xout{and} \over{a}{A}s the notion it-%
self is visibly a basic one,  there should be no question whatever to
change \deleted{it,} \inserted{the name} and replace it say by ``weakly aspheric'',  whereas the notion
is surely quite a strong one,  \deleted{which} \inserted{and} does'nt deserve such minimizing qua-%
lification !  There is indeed a stronger notion,  which in the context
of topological spaces or étale cohomology of schemes reduces to the pre-%
vious one\xout{, \unsure{par}} in the \inserted{particular} case of a map which is supposed \uline{proper} .  This con-%
dition could be expressed by saying that for \uline{any} base change $Y'\hand{\to}Y$

\margin{94}
